\subsection{Interrupting a Reflective Control Form}
\label{sec:interrupt}

The second experiment in \filepath{examples/misc/reify-args.abclr2}
defines \code{with-interrupt}. Unlike the preceding \code{my-if}, it
customizes an individual object's metaobject rather than the group's
evaluator. It illustrates \emph{abandoning} an evaluation and selecting
an alternative result. It is not an operating-system interrupt handler,
an exception mechanism with stack unwinding, or state rollback.

\subsubsection{Read the Application Before the Machinery}

The original worker counts, then returns one of two strings:

\wholelisting[firstline=194,lastline=202]{examples/misc/reify-args.abclr2}

The three arguments are the body, the success expression, and the
failure expression. They are passed as compiled representations, not
evaluated argument values. The body prints \code{Count 0.} through
\code{Count n-1.} if allowed to finish. Its ordinary \code{dotimes}
result is not the application's answer: the selected branch supplies
the value of \code{with-interrupt}, and the outer \code{!} sends that
value to the waiting caller.

The original controller starts the worker and later interrupts it:

\wholelisting[firstline=204,lastline=212]{examples/misc/reify-args.abclr2}

The worker's request is asynchronous with \code{@ R}, so the worker
answers the original caller directly. The controller does not wait for
the worker before running its dummy loop. After \code{m} empty iterations,
it addresses \code{[meta test-with-interrupt]} directly. Sending an
ordinary interrupt request to the application would queue it behind
the running script; the direct meta-level request can change evaluation
policy while that script is active. It has no application reply of its
own, and the controller does not execute a second \code{!}.

The commented historical session demonstrates these two outcomes:

\begin{codeblock}
\begin{lstlisting}
[test-with-interrupt-2 <== [:run 10 3]]
; Count 0., Count 1., Count 2.
; Being interrupted... / Context is caught.
; => "Time is up."

[test-with-interrupt-2 <== [:run 10 30]]
; Count 0. through Count 9.
; Interruption is ignored.
; => "Yeah! I did it!"
\end{lstlisting}
\end{codeblock}

The checked excerpt reproduces those counts and branch results on the
current image. The parameter \code{m} is a workload measured in empty
loop iterations, not seconds or a timeout deadline. Do not generalize
the exact point of interruption to another evaluator or scheduler.
Even diagnostic output order needs care: in the current transcript the
late interruption's ``ignored'' message appears after the check has
already received the success result. The worker's reply does not join
every outstanding action in the controller.

\subsubsection{How the Control Form Reaches the Metaobject}

The group's \code{meta-gen} selects \code{meta-with-interrupt-gen}.
The generator itself belongs to \code{[group-of group]}, the surrounding
group's meta-group. For each application object it constructs a
non-reifying evaluator wrapper and a non-reifying metaobject. These
machinery objects use direct message handlers, while the application
script can use reflective forms.
The macro exposes the protocol:

\wholelisting[firstline=190,lastline=192]{examples/misc/reify-args.abclr2}

Its expansion sends
\code{[:with-interrupt [Body Success Fail] Env Id Gid Eval] @ C}
to the current application's metaobject. \code{Env} is the captured
evaluation environment; \code{Id} identifies that application's
metaobject; \code{Gid} is its group context; \code{Eval} is its current
evaluator wrapper. \code{C} receives the result of the control form.
In this worker, that continuation leads to the outer \code{!}; it is
not simply the controller's application reply destination \code{R}.

Read the metaobject's control handlers in this order:

\wholelisting[firstline=105,lastline=121]{examples/misc/reify-args.abclr2}

On entry to \code{:with-interrupt}, it stores the failure representation
and the result/context tuple \code{[C Env Id Gid Eval]}. It requests
evaluation of \code{Body}; the attached continuation receives
\code{Result}, ignores that body result, and requests evaluation of
\code{Success} with destination \code{C}. The caller waits for one
selected branch, not for all three arguments.

An active metaobject accepts \code{[:interrupt]} and tells its evaluator
wrapper where to report an intercepted request. A dormant metaobject
only prints that interruption is ignored. The mode comes from the
ordinary message-dispatch lifecycle, not from a timer or a stack of
active \code{with-interrupt} regions.

\subsubsection{Capture the Next Request, Then Abandon It}

Here is the wrapper inside the original generator:

\wholelisting[firstline=78,lastline=93]{examples/misc/reify-args.abclr2}

The final delimiters also close the enclosing temporary-variable binding.
Initially \code{interrupted} is \code{nil}. A normal request \code{any}
is forwarded to the underlying group evaluator, retaining its reply
destination \code{R}. Here the capitalized \code{Eval} in the wrapper's
fallback comes from the generator's creation request: it names the
underlying evaluator. It is different from the incoming control
handler's \code{Eval}, which names the wrapper through which the
application is currently evaluated.

The \code{[:interrupt R]} handler saves its notification address and sets
a flag. It does not immediately stop host Lisp code or return a context.
The next ordinary evaluation request consumes the flag. Instead of
forwarding that request, the wrapper sends \code{[:context [any . R]]}
to the saved address. The dotted pair retains the request and its
completion destination. Interruption occurs at an evaluation-message
boundary; a host operation already executing is not preempted inside
its machine instructions.

The metaobject's \code{[:context \_]} handler intentionally ignores that
intercepted request. It evaluates \code{fail-form} using the separately
saved \code{context-for-recover}, with the same \code{C} as the control
form's original result destination. This bypasses the body continuation
that would have selected success. The interrupted computation is not
resumed later. Contrast this with the newer Express library, which
stores a suspended evaluation request and resumes it after the express
script has finished (Section~\ref{sec:original-express}).

Follow an early interruption through five distinct messages:

\begin{enumerate}
  \item The controller sends the worker \code{[:run n] @ R}.
  \item The worker's reflective form asks its metaobject to evaluate the
  body and remember a failure alternative.
  \item The controller sends that metaobject \code{[:interrupt]}; the
  metaobject sends its wrapper \code{[:interrupt Me]}.
  \item The wrapper intercepts a later evaluation request and reports
  \code{:context}; the metaobject substitutes failure evaluation at
  the control form's saved destination \code{C}.
  \item Failure returns through the worker's \code{!} to \code{R}.
  The script's separate execution-completion continuation eventually
  sends \code{:end} to its metaobject, advancing the ordinary queue.
\end{enumerate}

The remaining generator handlers reproduce the familiar queue protocol:
\code{:message} enqueues, \code{:begin} matches and starts one script,
and \code{:end} starts another request or returns to dormant mode.
State initialization uses \code{:do-evlis}, \code{pairlis}, and
\code{env-gen}, as in the inheritance walkthrough. Extra creation
parameters are matched by \code{Extra} but not used here. The
\code{:evaluator}/\code{:set-evaluator} handlers expose the evaluator
choice; they do not themselves define interruption.

\subsubsection{Effects, Scope, and Safe Reproduction}

Printed counts remain printed. A supplemental worker using the original
generator increments persistent state in its body and returns either
\code{(:SUCCESS count)} or \code{(:FAILED count)}. Its early interrupt
leaves a positive count below the requested 50 increments. A subsequent
ordinary state query sees that same count, and a new successful run adds
two. Thus branch substitution abandons further body work but does not
undo earlier effects or automatically reinitialize the worker.

The implementation has one \code{fail-form} and one recovery tuple per
metaobject, not a stack of nested regions. Use it around one body interval
rather than as a general nested cancellation API. There is no timeout clock, cleanup
protocol, or guarantee of arbitrary resource release.

For a separate executable version of this experiment, run
\code{make -C tutorial interrupt-check}. Then load the prepared definitions
in a fresh ABCL/R2 session:

\begin{codeblock}
\begin{lstlisting}
(inc "build/interrupt-tests/interrupt-excerpt" :use-source t)
\end{lstlisting}
\end{codeblock}

The excerpt contains the interrupt generator, its group, the worker and
the controller. Thirteen checks exercise the two controller calls,
empty and normal bodies, reuse, and the supplemental state experiment.
Sources, translations, inputs and logs stay in
\filepath{build/interrupt-tests/}.

\paragraph{Check your understanding.}
Why not use \code{reify-args} here? The handler is on the application's
metaobject, not the group evaluator. Why does the failure path not send
the intercepted request back to the wrapper? It substitutes another
expression rather than resuming that request. Does \code{m = 3} mean
three seconds? No. Does interruption restore the count? No; prior
effects remain. Why do both success and failure use \code{@ C}? Both
must complete the same control expression and resume its caller's
remaining evaluation, including the application's reply.
